\origin{ai}

\subsection{Forced period-ten clock without a canonical origin}
\label{sec:window6-clock-coupling-cert-basepoint-free-gauge}

\paragraph{Finite carrier.}
Let $\sigma$ be the cyclic left shift on binary words of length $10$. The ambient finite word space is $\{0,1\}^{10}$, and the cyclic golden-word carrier is the subset
$$
\begin{aligned}
P_{10}=\{w\in \{0,1\}^{10}:&\text{ no two adjacent cyclic coordinates of }w\text{ are both }1,\\
&\text{ and the }\sigma\text{-orbit of }w\text{ has period dividing }10\}.
\end{aligned}
$$
The clock orbit used here is the single-one orbit
$$
\begin{aligned}
O=\{\sigma^j(0000000001):0\leq j<10\}\subset P_{10}.
\end{aligned}
$$
It has exactly $10$ elements. The action of the cyclic group $C_{10}$ on $O$ generated by $\sigma$ is simply transitive: for every ordered pair $(o,o')\in O\times O$ there is a unique residue $j\in \{0,1,\ldots,9\}$ such that $o'=\sigma^j(o)$. Thus $O$ is a torsor for $C_{10}$, not a group with a preferred identity.

\paragraph{Exhaustive finite counts.}
The complete cyclic golden-word carrier has the exact size
$$
\begin{aligned}
|P_{10}|=123.
\end{aligned}
$$
The shift action on $P_{10}$ decomposes into orbits with the following period profile:
$$
\begin{aligned}
\{10:11,\;5:2,\;2:1,\;1:1\}.
\end{aligned}
$$
Equivalently, there are eleven orbits of period $10$, two orbits of period $5$, one orbit of period $2$, and one fixed orbit of period $1$. The orbit $O$ is one of the period-$10$ orbits, singled out by the predicate of having exactly one coordinate equal to $1$.

\paragraph{Clock matrix and spectral data.}
Choose any ordering of the ten elements of $O$ compatible with repeated application of $\sigma$. In that basis, let $S_{10}$ be the permutation matrix that sends the basis vector of $o$ to the basis vector of $\sigma(o)$. The matrix is a ten-cycle permutation matrix. Its characteristic polynomial is exactly
$$
\begin{aligned}
(x-1)(x+1)(x^4-x^3+x^2-x+1)(x^4+x^3+x^2+x+1).
\end{aligned}
$$
This spectral description records the cycle structure of the clock action. It gives eigenlines for the period-ten shift and includes the constant Perron vector, but it does not choose a distinguished element of $O$. The constant vector is invariant under changing the origin of the torsor, and the nonconstant eigenlines record phases only after an origin has already been selected.

\paragraph{Arithmetic register.}
The finite arithmetic register considered alongside the clock is the Fibonacci pair $(F_n,F_{n+1})$ modulo $10$. Its Pisano period is
$$
\begin{aligned}
\mathrm{Pisano}(10)=60.
\end{aligned}
$$
This gives a canonical period for the register sequence modulo $10$. It does not, by itself, define a map from the torsor $O$ to an absolute phase of that register. In particular, the existence of a period-$60$ register and a period-$10$ clock does not select one of the ten clock states as a zero state.

\paragraph{Candidate basepoint tests.}
Four natural finite candidates for a forward-distinguished basepoint are tested within this structure. First, fixed-point symmetries do not yield a unique point of $O$, because the simply transitive shift action makes every clock state structurally equivalent under the clock relation. Second, the torsor $O$ becomes a group only after an origin is chosen; the clock data do not supply such an identity element internally. Third, the modulo-$10$ Fibonacci register supplies a canonical periodic sequence but no intrinsic alignment map from $O$ to an absolute register phase. Fourth, the spectral clock supplies eigenlines and the constant Perron vector, but these data determine no canonical phase origin in the underlying ten-element orbit.

The lexicographic representative of $O$ is $0000000001$. If an external lexicographic order on ambient binary strings is imported, then the fixed seam word $\sigma^7(0000000001)$ has first-seam residue $7$. This is a valid calculation only after the ambient string order has been added. It is therefore not a forward-internal basepoint choice of the clock torsor itself.

\paragraph{Finite witness extraction.}
The certificate is an exhaustive finite procedure over the sets just described. It enumerates $P_{10}$ inside $\{0,1\}^{10}$ by the cyclic no-adjacent-$11$ predicate, partitions the result into $\sigma$-orbits, identifies the single-one orbit $O$, constructs the permutation matrix $S_{10}$ on $O$, computes the displayed characteristic polynomial by exact symbolic linear algebra, and checks the four basepoint candidates against the requirement that a basepoint be distinguished without importing an origin. No physical readout value, reverse fit, or metrological comparison is evaluated in this extraction.

\paragraph{Conclusion.}
The finite forced-window data determine a period-ten clock torsor, its shift partition inside $P_{10}$, the exact ten-cycle matrix spectrum, and the independent modulo-$10$ Fibonacci register period. They do not determine an absolute period-ten basepoint. The clock origin required by a clock-coupling certificate remains an irreducibly free gauge choice.

\paragraph{Local boundary.}
No identification with a physical constant, no global flow law, no external dependency, no metrological comparison, and no interpretation beyond this finite construction follows from the preceding calculation. In particular, a distinguished origin for the clock may be supplied only by a separate certificate that adds data not present in the finite clock torsor, the modulo-$10$ register, or the displayed spectrum.
